\section{A normalized-covariance proof of the nonlinear entropy bound}
\label{b:section}

We give a direct proof of the entropy-dependent estimate using the
supported-face codimension bound from
Proposition~\ref{a:prop-supported-codimension}.  The normalization below
makes the marginal-entropy Hessian equal to minus the identity, while
the covariance trace has an exact combinatorial value.  These two
identities lead to a contradiction at any positive maximum of the
proposed entropy defect.  The scalar deduction in
Corollary~\ref{cons:legacy-defect} is a separate proof of the same bound.

Throughout this appendix, \(G=(V,E)\) is a finite loopless multigraph
on \(2m>0\) vertices with at least one perfect matching.  Parallel
edges have distinct labels.  Write \(\mathcal M\) for its perfect
matchings, \(a(M)=\mathbf1_M\), and
\(P=P(G)=\operatorname{conv}\{a(M):M\in\mathcal M\}\).
Recall that, with natural logarithms and \(0\ln0=0\),
\[
 F(y)=-\sum_{e\in E}y_e\ln y_e,
 \qquad
 H(y)=\max_{\substack{p\text{ a probability law on }\mathcal M\\
                       \sum_Mp_Ma(M)=y}}
                   \left(-\sum_{M\in\mathcal M}p_M\ln p_M\right).
\]
To distinguish entropy of a law from its maximum at a prescribed mean,
write \(\operatorname{Ent}(p)=-\sum_Mp_M\ln p_M\).

\begin{theorem}[Nonlinear entropy bound]
\label{b:thm-nonlinear}
For every \(y\in P(G)\),
\begin{equation}
 F(y)-H(y)\le 8m\bigl(1-e^{-F(y)/m}\bigr).
 \label{b:eq-nonlinear}
\end{equation}
\end{theorem}

The attainment and continuity of \(H\) on the whole polytope follow
from Lemma~\ref{ent:continuity}. In particular, a continuous defect
has a maximum even when that maximum lies on a proper face. The proof
below combines the codimension bound with that lemma, the face
description, and the exponential law of Lemma~\ref{ent:exponential}.
These inputs are independent of Theorem~\ref{ent:main}.

\subsection{The face at a hypothetical positive maximum}
\label{b:positive-maximum}

\begin{proof}[Proof of Theorem~\ref{b:thm-nonlinear}]
Define
\[
 f(t)=8(1-e^{-t}),
 \qquad
 \Psi(y)=F(y)-H(y)-mf\bigl(F(y)/m\bigr).
\]
Suppose that \(\Psi\) is positive somewhere.  By
Lemma~\ref{ent:continuity}, it has a positive maximum at some
\(x\in P\).  Let \(F_x\) be the minimal face containing \(x\), and put
\[
 A=\{e\in E:x_e>0\},\qquad s=|A|,\qquad
 d=\dim F_x,\qquad t=F(x)/m.
\]
Here \(s\) counts supported \emph{edges}.  Each matching has exactly
\(m\) edges, so \(\sum_ex_e=m\) and \(0\le x_e\le1\).
Concavity of \(-u\ln u\) gives
\[
 0\le t\le\ln(s/m).
\]
On \([0,\ln8]\), we have \(f(0)=0\) and \(f'(t)\ge1\), so
\(f(t)\ge t\).  Since \(H(x)\ge0\), positivity of \(\Psi(x)\)
therefore forces \(t>\ln8\), and hence \(s>8m\).
The supported-face estimate of
Proposition~\ref{a:prop-supported-codimension} now yields
\begin{equation}
 d\ge s-4m>\frac{s}{2}>0.
 \label{b:eq-large-dimension}
\end{equation}
Thus all integral and low-entropy cases have already been excluded
before any covariance inverse is introduced.

Let
\[
 \mathcal M_x=\{M\in\mathcal M:a(M)\in F_x\}.
\]
Every law whose mean is in \(F_x\) is supported on \(\mathcal M_x\).
Indeed, each inequality active on this face has nonnegative slack at
every matching vector.  Expected slack zero forces every atom of
positive probability to satisfy its equality.  This applies to both
zero coordinates and tight odd cuts.  Consequently
\(F_x=\operatorname{conv}\{a(M):M\in\mathcal M_x\}\).

Discard the coordinates outside \(A\) and set
\begin{equation}
 D=\operatorname{diag}(x_e:e\in A),\qquad
 \mathcal T=D^{-1/2}T(x)
     =\operatorname{span}\bigl(D^{-1/2}(F_x-x)\bigr)
       \subseteq\mathbb R^A.
 \label{b:eq-normalization}
\end{equation}
All diagonal entries of \(D\) are positive.  By
Lemma~\ref{geom:lem-face}, \(x\) is relatively interior to \(F_x\).
Therefore
\[
 z=D^{-1/2}(y-x),\qquad y=x+D^{1/2}z
\]
identifies a relative neighborhood of \(x\) in the face with a
neighborhood of zero in the \(d\)-dimensional Euclidean space
\(\mathcal T\).  Define
\[
 b_M=D^{-1/2}(a(M)-x),\qquad M\in\mathcal M_x.
\]
Their convex hull contains a neighborhood of zero in \(\mathcal T\).
All gradients, Hessians, operator inverses, and operator traces in
the remainder of the proof are taken on \(\mathcal T\).
Coordinates with \(x_e=1\) have not been discarded: their centered
entries and all tangent directions are zero, so they cause no
singularity on this space.

\subsection{A finite exponential family on the normalized tangent space}
\label{b:exponential-family}

We now express the face law of Lemma~\ref{ent:exponential} in the
normalized coordinates. Let \(p\) be the positive maximizing law at
\(x\), and write its parameter as \(\theta\in\mathbb R^A\),
restricting to supported coordinates. Since
\(a(M)=x+D^{1/2}b_M\),
\[
 \theta\cdot a(M)=\theta\cdot x+
                      (D^{1/2}\theta)\cdot b_M.
\]
Every \(b_M\) lies in \(\mathcal T\). Thus only the orthogonal
projection \(\lambda=\operatorname{proj}_{\mathcal T}
(D^{1/2}\theta)\) affects the normalized law. Define
\[
 \Phi(\eta)=\ln\sum_{M\in\mathcal M_x}e^{\eta\cdot b_M}
 \qquad(\eta\in\mathcal T).
\]
The law and its normalized mean then satisfy
\begin{equation}
 p_M=\exp\bigl(\lambda\cdot b_M-\Phi(\lambda)\bigr),
 \qquad \nabla\Phi(\lambda)=0.
 \label{b:eq-exponential-law}
\end{equation}
The last equality is the zero mean of \(b_M\) under \(p\).
The parameter is finite because the original law is taken on its
minimal face.

For any \(\eta\in\mathcal T\), put
\[
 p_{\eta,M}=e^{\eta\cdot b_M-\Phi(\eta)},\qquad
 z(\eta)=\sum_Mp_{\eta,M}b_M,\qquad
 y(\eta)=x+D^{1/2}z(\eta).
\]
This law has exponential parameter \(D^{-1/2}\eta\) in the
original coordinates, since the centering term is constant across
atoms. Lemma~\ref{ent:exponential} therefore shows that it maximizes
entropy at its own mean \(y(\eta)\) and remains on the same minimal
face. Differentiating the finite sum defining \(\Phi\) yields
\begin{equation}
 \begin{aligned}
 z(\eta)&=\nabla\Phi(\eta),\\
 H(y(\eta))&=\Phi(\eta)-\eta\cdot z(\eta),\\
 C(\eta)&:=\nabla^2\Phi(\eta)
       =\operatorname{Cov}_{p_\eta}(b_M).
 \end{aligned}
 \label{b:eq-mean-entropy-covariance}
\end{equation}
Each \(C(\eta)\) is positive definite on \(\mathcal T\).
Indeed, zero variance in a direction \(u\in\mathcal T\) would make
\(u\cdot b_M\) constant on all face atoms because \(p_\eta\) is
positive.  The affine span of the \(b_M\) is \(\mathcal T\), so
this forces \(u=0\).  Write \(C=C(\lambda)\).

\subsection{The second variation and its trace}
\label{b:second-variation}

We have a positive definite covariance on a space of dimension
\(d>s/2\).  We now compare its inverse with the marginal-entropy
Hessian.  This comparison uses the smooth exponential parameter;
it does not require differentiability of \(H\) across faces.

Set
\[
 \widehat F(z)=F(x+D^{1/2}z),\qquad
 L(z)=\widehat F(z)-mf\bigl(\widehat F(z)/m\bigr).
\]
These functions are smooth near zero on \(\mathcal T\).  In the
retained coordinates, the Hessian of \(F\) at \(x\) is \(-D^{-1}\).
The change of coordinates in Equation~\eqref{b:eq-normalization}
therefore gives
\begin{equation}
 \nabla^2\widehat F(0)=-I_{\mathcal T}.
 \label{b:eq-normalized-hessian}
\end{equation}
By Equation~\eqref{b:eq-mean-entropy-covariance}, the smooth function
\[
 J(\eta)=L(z(\eta))-\Phi(\eta)+\eta\cdot z(\eta)
        =\Psi(y(\eta))
\]
has a local maximum at \(\lambda\).  Its gradient is
\[
 \nabla J(\eta)
   =C(\eta)\bigl(\nabla L(z(\eta))+\eta\bigr).
\]
Since \(C\) is invertible and \(z(\lambda)=0\), stationarity implies
\(\nabla L(0)+\lambda=0\).  Differentiating once more at
\(\lambda\), the derivative of \(C(\eta)\) is multiplied by this
zero vector.  Thus the second-derivative test gives
\[
 \nabla^2J(\lambda)=C\nabla^2L(0)C+C\preceq0.
\]
Congruence by \(C^{-1}\), as an operator on \(\mathcal T\), yields
\begin{equation}
 \nabla^2L(0)+C^{-1}\preceq0.
 \label{b:eq-hessian-inequality}
\end{equation}

Under \(p=p_\lambda\), the normalized vector \(b_M\) has mean zero.
For each edge \(e\), its indicator in the random matching has variance
\(x_e(1-x_e)\), so
\begin{equation}
 \operatorname{tr}_{\mathcal T}C
 =\mathbb E_p\|b_M\|_2^2
 =\sum_{e\in A}\frac{x_e(1-x_e)}{x_e}
 =s-m.
 \label{b:eq-covariance-trace}
\end{equation}
The trace on \(\mathcal T\) equals the ambient trace because all
the vectors \(b_M\) lie in \(\mathcal T\).
For \(u>0\), the inequality \(u^{-1}\ge2-u\) follows from
\((u-1)^2/u\ge0\).  Applying it to the \(d\) positive eigenvalues
of \(C\) gives
\begin{equation}
 \begin{split}
 \operatorname{tr}_{\mathcal T}(-I_{\mathcal T}+C^{-1})
 &\ge d-\operatorname{tr}_{\mathcal T}C\\
 &=d-s+m\ge-3m.
 \end{split}
 \label{b:eq-inverse-trace}
\end{equation}
The last inequality is again the supported-face codimension bound.

It remains to include the nonlinear correction.  Writing
\(g=\nabla\widehat F(0)\), the full chain rule is
\begin{equation}
 \left.\nabla^2\!\left[mf\bigl(\widehat F(z)/m\bigr)\right]
       \right|_{z=0}
 =-f'(t)I_{\mathcal T}+\frac{f''(t)}{m}gg^{\mathsf T}.
 \label{b:eq-rank-one-correction}
\end{equation}
In particular, the rank-one term has trace
\(f''(t)\|g\|_2^2/m\le0\).  Since
\(t\le\ln(s/m)\), we also have \(f'(t)=8e^{-t}\ge8m/s\).
Combining Equations~\eqref{b:eq-large-dimension},
\eqref{b:eq-inverse-trace}, and
\eqref{b:eq-rank-one-correction}, we obtain
\begin{equation}
 \begin{split}
 \operatorname{tr}_{\mathcal T}\bigl(\nabla^2L(0)+C^{-1}\bigr)
 &\ge-3m+f'(t)d\\
 &\ge-3m+\frac{8md}{s}>m>0.
 \end{split}
 \label{b:eq-strict-contradiction}
\end{equation}
This contradicts the negative semidefiniteness in
Equation~\eqref{b:eq-hessian-inequality}.  Hence \(\Psi(y)\le0\)
throughout \(P\), proving Theorem~\ref{b:thm-nonlinear}.
\end{proof}

The theorem assumes \(m>0\) because its formula divides by \(m\).
On the empty vertex set, the unique empty matching has
\(F=H=0\); its entropy statement is handled separately without that
quotient.  If a nonempty graph has no perfect matching, there are no
means in \(P(G)\) to which the theorem applies.
